\section{Síntesis y ampliación}

La semana 7 exploró la evolución del papel de la revisión de código en la era de la codificación con IA. El mensaje central es: la IA hace la revisión de código más eficiente, pero al mismo tiempo la vuelve más indispensable, porque hay que revisar no solo el código escrito por personas, sino también el código escrito por la IA.

\subsection{Puntos clave}

\begin{itemize}
    \item La tasa de detección de errores de la revisión de código (55--60\%) supera con mucho a la de las pruebas (25--45\%)
    \item Un buen comentario de revisión es constructivo, específico e invita al diálogo
    \item La revisión debe cubrir cinco dimensiones: corrección, legibilidad, rendimiento, seguridad y mejores prácticas
    \item La revisión de código con IA es adecuada para atender problemas «superficiales»; las personas se concentran en problemas «profundos»
    \item Limitaciones actuales: falsos positivos, incapacidad de captar idiomas propios del repositorio y de manejar lógica de negocio compleja
    \item En la era de la IA la revisión de código es más importante; «lo escribió la IA» no es una excusa que exima de responsabilidad
\end{itemize}

\subsection{Recomendaciones para la implementación en equipo}

\begin{enumerate}
    \item Primero estandarizar la plantilla de PR, la descripción de cambios y la descripción de pruebas, y solo entonces introducir la revisión con IA
    \item Dejar que la IA cubra primero los problemas superficiales de alta frecuencia; no exigirle desde el inicio que sustituya al revisor experimentado
    \item Construir una cadena de evidencia para los comentarios clave: citar el diff, la regla o el contexto histórico
    \item Medir de manera continua la tasa de falsos positivos y de omisiones, y tratar el sistema de revisión como un producto que se opera
    \item Mantener un umbral de aprobación humana más estricto en módulos de alto riesgo como seguridad, pagos y permisos
\end{enumerate}

\subsection{Tabla comparativa de estrategias de revisión}

\begin{table}[H]
\centering
\begin{tabular}{p{0.24\textwidth} p{0.27\textwidth} p{0.31\textwidth}}
\toprule
Escenario & Tareas más adecuadas para la IA & Tareas más adecuadas para las personas \\
\midrule
PR de negocio habituales & Nomenclatura, lógica duplicada, errores comunes, verificación de reglas & Intención de negocio, decisiones de límites, dirección de evolución del código \\
Cambios de infraestructura & Validación de configuración, análisis del impacto en dependencias & Estrategia de reversión, riesgo de compatibilidad, ritmo de la migración \\
Módulos sensibles a la seguridad & Coincidencia de patrones de vulnerabilidades conocidas, aviso de cambios de permisos & Modelado de amenazas, evaluación de la superficie de ataque, aprobación de excepciones \\
Refactorización a gran escala & Verificación de consistencia de patrones, detección de referencias faltantes & Si vale la pena refactorizar, estrategia de división, ventana de despliegue \\
\bottomrule
\end{tabular}
\caption{Frontera de la división de trabajo entre el reviewer de IA y el reviewer humano}
\end{table}

\subsection{Lecturas adicionales}

\begin{itemize}
    \item Graphite: \url{https://graphite.dev}
    \item Greptile: \url{https://greptile.com}
    \item CodeRabbit: \url{https://coderabbit.ai}
    \item Prácticas recomendadas de revisión de código de Google: \url{https://google.github.io/eng-practices/review/}
    \item ``How to do a code review'' by Blake Smith
\end{itemize}

%% --- Fin del contenido --- %%

\end{document}
